% ─────────────────────────────────────────────────────────────────────────────
%  Apéndice A – Gramática BNF completa del subconjunto de Java
% ─────────────────────────────────────────────────────────────────────────────
\chapter{Gramática BNF del subconjunto Java}
\label{apend:gramatica}

Las reglas se presentan en notación EBNF. Los nombres entre
\texttt{<>} son no terminales; los literales entre comillas simples
son terminales. Los corchetes denotan opcionalidad; las llaves,
cero o más repeticiones; las alternativas se separan con \texttt{|}.

\section{Estructura del programa}

\begin{lstlisting}[style=pseudo, caption={Producción raíz}]
<Program>  ::= 'public' 'class' ID '{' <Members>* '}'

<Members>  ::= <Annotation>
            |  <Method>

<Annotation> ::= '@pre'       <Expr>
              |  '@post'      <Expr>
              |  '@invariant' <Expr>
              |  '@variant'   <Expr>
\end{lstlisting}

\section{Declaración de métodos}

\begin{lstlisting}[style=pseudo, caption={Métodos y parámetros}]
<Method>  ::= 'public' 'static' <Type> ID '(' [<Params>] ')' <Block>
           |  'public' 'static' 'void' ID '(' [<Params>] ')' <Block>

<Params>  ::= <Type> ID { ',' <Type> ID }

<Type>    ::= 'int' | 'double' | 'boolean' | 'String[]'
\end{lstlisting}

\section{Sentencias}

\begin{lstlisting}[style=pseudo, caption={Sentencias del subconjunto}]
<Block>   ::= '{' <Stmt>* '}'

<Stmt>    ::= <VarDecl>
           |  <Assign>
           |  <If>
           |  <While>
           |  <Return>
           |  <MethodCall> ';'
           |  <PrintStmt>

<VarDecl> ::= <Type> ID ';'
<Assign>  ::= ID '=' <Expr> ';'

<If>      ::= 'if' '(' <Expr> ')' <Block>
           |  'if' '(' <Expr> ')' <Block> 'else' <Block>

<While>   ::= 'while' '(' <Expr> ')' <Annotation>* <Block>

<Return>  ::= 'return' <Expr> ';'

<PrintStmt> ::= 'System.out.println' '(' <Expr> ')' ';'
             |  'System.out.print'   '(' <Expr> ')' ';'
\end{lstlisting}

\section{Expresiones}

\begin{lstlisting}[style=pseudo, caption={Gramática de expresiones (precedencia ascendente)}]
<Expr>  ::= <Expr> '||' <Expr>
         |  <Expr> '&&' <Expr>
         |  '!' <Expr>
         |  <Expr> ('==' | '!=' | '<' | '<=' | '>' | '>=') <Expr>
         |  <Expr> ('+' | '-') <Expr>
         |  <Expr> ('*' | '/' | '%') <Expr>
         |  '-' <Expr>             -- negacion unaria
         |  '(' <Expr> ')'
         |  ID
         |  INT_LIT
         |  REAL_LIT
         |  'true' | 'false'
         |  'result'               -- solo en anotaciones @post
         |  <MethodCall>

<MethodCall> ::= ID '(' [<Args>] ')'
<Args>       ::= <Expr> { ',' <Expr> }
\end{lstlisting}

\section{Anotaciones de verificación}

Las anotaciones de verificación (Milestone~5) son un mecanismo de
extensión léxica: los tokens \texttt{@pre}, \texttt{@post},
\texttt{@invariant} y \texttt{@variant} se reconocen por el analizador
léxico (Flex) y se integran en la gramática como sentencias de primer
nivel dentro de \texttt{Members}, o como prefijos de un \texttt{While}.

Las expresiones de anotación comparten la misma gramática que
\texttt{<Expr>}, con el añadido del símbolo \texttt{result} que
referencia el valor de retorno del método en la postcondición.

\section{Precedencia y asociatividad de operadores}

\begin{table}[H]
  \centering
  \caption{Precedencia (mayor nivel = mayor prioridad)}
  \label{tab:prec}
  \begin{tabular}{cll}
    \toprule
    \textbf{Nivel} & \textbf{Operadores} & \textbf{Asociatividad} \\
    \midrule
    1 & \texttt{||} & izquierda \\
    2 & \texttt{\&\&} & izquierda \\
    3 & \texttt{!} & derecha (prefijo) \\
    4 & \texttt{== !=} & izquierda \\
    5 & \texttt{< <= > >=} & izquierda \\
    6 & \texttt{+ -} & izquierda \\
    7 & \texttt{* / \%} & izquierda \\
    8 & unario \texttt{-} & derecha (prefijo) \\
    9 & llamada a función, literal, identificador & --- \\
    \bottomrule
  \end{tabular}
\end{table}
